\chapter{TRABAJO FUTURO} \label{TrabajoFuturo}

En este trabajo principalmente se ha realizado una investigación para definir la aplicación de apoyo en la búsqueda y rescate de personas desaparecidas y un prototipo básico, por lo que todavía existe recorrido para conseguir una solución realmente fiable y operativa.

El primer paso consiste en consolidar el prototipo actual añadiendo las funcionalidades pendientes como el inicio de sesión, la diferenciación de perfiles, la geolocalización de participantes en una búsqueda, las notificaciones push o la posibilidad de crear puntos de control en la búsqueda, todo ello implicará crear las tablas necesarias en la base de datos AlertRes de MariaDB y adaptar el frontend y el backend de la aplicación. Una vez obtenido un prototipo suficientemente robusto y funcional se deberán realizar pruebas de uso con equipos reales para evaluar la usabilidad, rendimiento, recoger feedback y corregir errores.

Paralelamente, se deberá generar conjuntos de datos simulados para comenzar con el estudio del sistema de predicción (descrito más adelante en este capítulo) de la localización de una persona desaparecida en base a las características del caso, de la propia persona y datos históricos anónimos.

Cuando se encuentre desarrollado completamente el prototipo, se redefinirá el entorno de desarrollo, las plataformas soportadas y la compatibilidad con APIs externas y otros sistemas de datos. Además, para garantizar respaldo institucional y el acceso a datos reales que ayuden con el desarrollo del modelo de predicción, conviene buscar acuerdos con organismos competentes como el CNDES.

Finalmente, una vez alcanzada una madurez técnica e institucional de la aplicación, se pueden contemplar la incorporación de nuevas funcionalidades ya definidas que aumenten la eficacia operativa, o algunas nuevas no definidas previamente como algoritmos que sugieran la estrategia de búsqueda óptima en función de las características del caso, integración de dispositivos de monitorización de participantes y recursos o diseño de sensores para perros de búsqueda cuya información quede registrada y asociada en la aplicación.

\section{Propuesta de sistema de predicción}

Aunque actualmente se carece de acceso a datos reales y a bases de datos policiales como PDyRH, se plantea un método de predicción de forma teórica que podría implementarse en futuras versiones de la aplicación AlertRes. Este método se fundamenta en el análisis de datos propios de la aplicación (casos registrados, perfiles de desaparecidos, condiciones ambientales, zonas de búsqueda, etc.) y, potencialmente, en la integración de fuentes externas como bases de datos institucionales, siempre que se cumplan los requisitos legales y éticos de acceso y protección de datos.

Se propone un enfoque basado en algoritmos de machine learning, entrenados con datos del histórico anonimizado de la aplicación o utilizando datos anonimizados de bases de datos policiales, lo que requeriría de algún tipo de convenio con el CNDES. El objetivo sería estimar el grado de urgencia del caso y obtener zonas con mayor probabilidad de localización para poder priorizar recursos y áreas de búsqueda.

Como se ha mencionado previamente para este estudio existen limitaciones importantes, puesto que la aplicación aún no es operativa, no se dispone de datos reales de bases de datos institucionales. Sin embargo, en el caso de llegar a disponer de datos, las variables de entrada incluirían, entre otras, datos sociodemográficos del desaparecido (edad, sexo...), información clínica relevante (patologías, movilidad, tratamiento...), características del incidente (última localización conocida, hora y fecha de desaparición, circunstancias reportadas...), atributos espaciales y ambientales (orografía, cobertura vegetal, redes hidrográficas, condiciones meteorológicas...), y metadatos operativos (tiempo desde la notificación, recursos movilizados, resultados de búsquedas previas...). Todas las variables sensibles se tratarían mediante anonimización para no exponer identidades. Las variables se irán incorporando de forma progresiva, realizando estudios previos para evaluar su contribución.

Será necesario definir una hoja de ruta del estudio que incluya la selección de qué datos son más interesantes para los inicios, realizar un preprocesamiento de los datos con división en conjuntos de datos de entrenamiento y validación, entrenar, validar, comparar y seleccionar los mejores modelos (Random Forest, KNN, SVM, redes neuronales...). También es interesante integrar el sistema predictivo con parámetros ya usados actualmente en las búsquedas, como Probabilidad de Área (POA) o Probabilidad de detección (POD), Probabilidad de éxito (POS)... \cite{ManualDeBusquedayRescate}

Sin embargo, conviene enfatizar que lo descrito en este apartado es una propuesta y no existen resultados empíricos que avalen la eficacia de un sistema predictivo de estas características. En caso de que se obtuvieran resultados exitosos, cualquier implementación real debería incorporar controles de gobernanza de datos, evaluación de sesgos, protocolos de anonimización robustos y pruebas piloto con supervisión de las autoridades competentes.